%----------------------------------------------------------------------------------------
%    PACKAGES AND THEMES
%----------------------------------------------------------------------------------------

\documentclass[aspectratio=169,xcolor=dvipsnames]{beamer}

\AtBeginSection[]
{
    \begin{frame}
        \frametitle{Visão Geral}
        \tableofcontents[currentsection]
    \end{frame}
}

\usefonttheme{professionalfonts} % using non standard fonts for beamer
\usefonttheme{serif} % default family is serif
\usepackage{fontspec}
\setmainfont{XCharter}[
    UprightFont    = *-Roman,
    BoldFont       = *-Bold,
    ItalicFont     = *-Italic,
    BoldItalicFont = *-BoldItalic,
]

\usepackage{hyperref}
\usepackage{graphicx} % Allows including images
\usepackage{booktabs} % Allows the use of \toprule, \midrule and \bottomrule in tables
\input{./slides/SimplePlusTheme/beamerthemeSimplePlus.sty}
\input{./slides/SimplePlusTheme/beamerinnerthemeSimplePlus.sty}
\input{./slides/SimplePlusTheme/beamerfontthemeSimplePlus.sty}
\input{./slides/SimplePlusTheme/beamercolorthemeSimplePlus.sty}

\usepackage{listings}
\setbeamercovered{transparent}

\usepackage{biblatex}
\addbibresource{slides/reference.bib}
\renewcommand*{\bibfont}{\footnotesize}
% Tamanho dos slides de "Leituras Recomendadas": maior que a bibliografia
% completa (\bibfont acima, que rende \tiny via a redefinição de \footnotesize)
% para destacar as leituras curadas.
\newcommand{\leiturasize}{\small}

\usepackage[brazilian ]{babel}
\usepackage{csquotes}
\usepackage{pgfplots}
\pgfplotsset{compat=1.18}

\usepackage{emoji}
\usepackage{amsmath}
\usepackage[xcharter,bigdelims,vvarbb]{newtxmath}
% newtxmath's operators font requests the fontspec family `XCharter(0)' in OT1;
% without these substitutions math digits and operator names silently fall back
% to Computer Modern (LaTeX Font Warning: Font shape `OT1/XCharter(0)/m/n' undefined).
\DeclareFontFamily{OT1}{XCharter(0)}{}
\DeclareFontShape{OT1}{XCharter(0)}{m}{n}{<->ssub * XCharter-TLF/m/n}{}
\DeclareFontShape{OT1}{XCharter(0)}{m}{it}{<->ssub * XCharter-TLF/m/it}{}
\DeclareFontShape{OT1}{XCharter(0)}{b}{n}{<->ssub * XCharter-TLF/b/n}{}
\DeclareFontShape{OT1}{XCharter(0)}{bx}{n}{<->ssub * XCharter-TLF/b/n}{}

\usepackage{bm}
\usepackage[table]{xcolor}
\usepackage{xcolor}
\usepackage{algpseudocode}

\usepackage{cancel}
\usepackage{ulem}

\usetikzlibrary{shapes}
\usetikzlibrary{arrows}
\usetikzlibrary {arrows.meta}
\usetikzlibrary{calc,positioning}
\usepgfplotslibrary{fillbetween}
\usetikzlibrary{angles,quotes}
\usetikzlibrary{tikzmark}


\definecolor{PastelRed}{HTML}{FFC1CC}
\definecolor{PastelBlue}{HTML}{C2F0FF}
\definecolor{PastelBlue2}{HTML}{3182BD}

\definecolor{PastelPink}{HTML}{FFCBE1}
\definecolor{PastelGreen}{HTML}{D6E5BD}
\definecolor{PastelYellow}{HTML}{F9E1A8}
\definecolor{PastelBlue}{HTML}{BCD8EC}
\definecolor{PastelPurple}{HTML}{DCCCEC}
\definecolor{PastelOrange}{HTML}{FFDAB4}


\renewcommand{\footnotesize}{\tiny}

\newcommand{\mespor}{
  \ifcase\month
    \or janeiro
    \or fevereiro
    \or março
    \or abril
    \or maio
    \or junho
    \or julho
    \or agosto
    \or setembro
    \or outubro
    \or novembro
    \or dezembro
  \fi
}
% \usepackage{csquotes}
% \usepackage{minted}
% \usemintedstyle{xcode}
\definecolor{vscLightBackground}{HTML}{FFFFFF}
\definecolor{vscLightText}{HTML}{000000}
\definecolor{vscLightGreen}{HTML}{008000}        % For Comments AND Numbers
\definecolor{vscLightRed}{HTML}{A31515}          % For Strings
\definecolor{vscLightTeal}{HTML}{267F99}
\definecolor{vscBlue}{HTML}{0000FF}% For main keywords (if, for, def...)
\definecolor{vscLightPytorch}{HTML}{800080}      % For PyTorch keywords (dark magenta)

\lstdefinestyle{vscode-light}{
    language=Python,
    backgroundcolor=\color{vscLightBackground},
    basicstyle=\ttfamily\scriptsize\color{vscLightText},
    keywordstyle=\color{vscLightTeal},
    commentstyle=\color{vscLightGreen},
    stringstyle=\color{vscLightRed},
    % Style for PyTorch keywords (using keyword group 2)
    keywordstyle=[2]\color{vscLightPytorch}, 
    morekeywords={self, True, False, None},
    morekeywords=[2]{
        torch, Tensor, device, dtype, to, nn, optim, autograd, utils, data,
        Module, Linear, Conv2d, ReLU, CrossEntropyLoss, Adam, SGD,
        Dataset, DataLoader, randn, zeros, ones, arange, cat, stack,
        view, reshape, sum, mean, max, min, matmul, no_grad, backward
    },
    keywordstyle=[3]\color{vscBlue}, 
    morekeywords=[3]{
        model, loss, y, y_hat, x , optimizer
    },
    frame=single,
    tabsize=1,
    showstringspaces=false,
    breaklines=true,
    captionpos=b,
    extendedchars=true,
    numbers=left,
}

\usepackage[cache=true]{minted}
% minted v3+ (TeX Live 2024+) auto-detects the output directory.
% minted v2 (Ubuntu apt TeX Live 2023) needs it set explicitly to
% match latexmkrc's $out_dir = 'build'.
\makeatletter
\@ifpackagelater{minted}{2024/01/01}{}{%
  \def\minted@outputdir{build/}%
}
\makeatother
\setminted{
    fontsize=\scriptsize,
    style=vs,
    breaklines=true,
    numbers=left,
    numbersep=5pt,
    frame=single,
}
\providecommand{\citeauthoryear}[1]{(\citeauthor{#1},~\citeyear{#1})}

\providecommand{\thetav}{\bm{\theta}}

% vetor coluna 2D com colchetes
\providecommand{\cvec}[2]{\begin{bmatrix} #1 \\ #2 \end{bmatrix}}

\providecommand{\varvector}[1]{\mathbf{#1}}
% instance
\providecommand{\X}{\varvector{X}}
\providecommand{\mlinstance}{\varvector{X}}
\providecommand{\mlinstancei}{\varvector{X}^{(i)}}
\providecommand{\mlinstancex}[1]{\varvector{X}^{(#1)}}

% target
\providecommand{\y}{\varvector{y}}
\providecommand{\mltarget}{\varvector{y}}
\providecommand{\mltargeti}{\varvector{y}^{(i)}}
\providecommand{\mltargetx}[1]{\varvector{y}^{(#1)}}

%hats
\providecommand{\yhat}{\hat{\varvector{y}}}
\providecommand{\mltargethat}{\hat{\varvector{y}}}
\providecommand{\mltargethati}{\hat{\varvector{y}}^{(i)}}
\providecommand{\mltargethatx}[1]{\hat{\varvector{y}}^{(#1)}}

% linear reg
\providecommand{\mllinearregression}{\mltargethat=\mlinstance\bm{\theta}}

\providecommand{\tikzarrowcolor}[3]{%
  \begin{tikzpicture}[remember picture, overlay]
    % Arrow body
    \path[fill=#3] ([shift={(#1,#2)}]current page.south west) 
      rectangle ++(0.3, -0.2);
    % Arrow head
    \path[fill=#3] ([shift={(#1+0.3,#2+0.1)}]current page.south west) -- 
         ++(0, -0.4) -- 
         ++(0.2, 0.2) -- 
         cycle;
  \end{tikzpicture}%
}

\providecommand{\tikzarrow}[2]{%
    \tikzarrowcolor{#1}{#2}{Red}
}



\DeclareMathOperator*{\argmax}{argmax}
\DeclareMathOperator*{\argmin}{argmin}




\providecommand{\mlexampleinstance}{ \underset{N\times M}{
    \begin{array}{|c|c|}
        \hline
        1 & 1   \\
        \hline
        1 & 2   \\
        \hline
        1 & 3   \\
        \hline
      \end{array}
      }}
      
\providecommand{\mlexampletarget}{      \underset{N\times 1}{
      \begin{array}{|c|}
        \hline
        20   \\
        \hline
        10   \\
        \hline
        30   \\
        \hline
      \end{array}}}
      
\providecommand{\mlexampleinstancetranspose}{\underset{M\times N}{
      \begin{array}{|c|c|c|}
        \hline
        1 & 1 & 1  \\
        \hline
        1 & 2 & 3\\
        \hline
      \end{array}}}



%----------------------------------------------------------------------------------------
%    TITLE PAGE
%----------------------------------------------------------------------------------------

\title{Deep Learning}
\subtitle{Revisão de \textit{Machine Learning}}

\author{Lucas Silveira Kupssinskü}

\institute
{
    Machine Learning Theory and Applications Laboratory \\
    PUCRS % Your institution for the title page
}
\date{\mespor de \the\year}
 % Date, can be changed to a custom date

%----------------------------------------------------------------------------------------
%    PRESENTATION SLIDES
%----------------------------------------------------------------------------------------

\begin{document}

\begin{frame}
    % Print the title page as the first slide
    \titlepage
\end{frame}

\begin{frame}{Overview}
    % Throughout your presentation, if you choose to use \section{} and \subsection{} commands, these will automatically be printed on this slide as an overview of your presentation
    \tableofcontents
\end{frame}

%------------------------------------------------
\section{Por que precisamos de Machine Learning?}
%------------------------------------------------

\begin{frame}{A importância de Machine Learning}

\begin{itemize}
    \item ``A breakthrough in machine learning would be worth ten Microsofts'' (Bill Gates)
    \item ``Machine learning is the next Internet'' (Tony Tether, Director, DARPA)
    \item ``Machine learning is going to result in a real revolution'' (Greg Papadopoulos, Former CTO, Sun)
    \item ``Machine learning is today’s discontinuity'' (Jerry Yang, Founder, Yahoo)
    \item ``Machine learning is the new electricity'' (Andrew NG, Stanford Professor and Coursera founder)
\end{itemize}
\end{frame}

\begin{frame}{\textit{Software} sem \textit{Machine Learning}}
    É uma sequência de instruções que:
    \begin{itemize}
        \item Recebe uma entrada
        \item Realiza um processamento
        \item Retorna uma saída
    \end{itemize}
\end{frame}
\begin{frame}{\textit{Software} sem \textit{Machine Learning}}{Exemplo}
    \begin{enumerate}
        \item Se você ou o seu oponente tiverem dois em sequência, jogue no espaço restante
        \item Caso contrário, se existir um movimento que crie duas sequencias de dois, faça essa jogada
        \item Caso contrário, se o centro estiver vazio, jogue no centro
        \item Caso contrário, se o seu oponente jogou em um canto, jogue no canto oposto 
        \item Caso contrário, se houver um canto livre jogue nesse canto
        \item Caso contrário, jogue em uma posição vazia qualquer
    \end{enumerate}
\end{frame}

\begin{frame}{Algumas vezes não sabemos como resolver}
\begin{columns}
    \column{0.5\textwidth}
       \begin{itemize}
            
           \item Complexidade de Tempo
           \item Complexidade de Espaço
           \item Complexidade ``Humana''
       \end{itemize}
    
    \column{0.5\textwidth}
        \begin{figure}
        \centering
         
\includegraphics[width=\textwidth]{slides/img/aula00/three-headed-dragon-meme.jpg}      
         \end{figure}
    
\end{columns}
\end{frame}

\begin{frame}{E no Aprendizado de Máquina?}

\begin{itemize}
    \item Mesmo quando não é possível criar um algoritmo para resolver o problema diretamente, pode ser possível criar algoritmos que \alert{\textit{encontram padrões}} 
    \item Usamos raciocínio \alert{indutivo} para resolver os problemas 
    
\end{itemize}

    \begin{block}{Definição de \textit{Machine Learning}}
        ``A computer program is said to learn from experience \textbf{\textcolor{red}{E}} with respect to some class of tasks \textbf{\textcolor{blue}{T}} and performance measure \textbf{\textcolor{green}{P}}, if its performance at tasks in \textbf{\textcolor{blue}{T}}, as measured by \textbf{\textcolor{green}{P}}, improves with experience \textbf{\textcolor{red}{E}}'' \citeauthoryear{mitchell_machine_1997}
    \end{block}
    
    
\end{frame}

\begin{frame}{Tipos de Raciocínio}
\begin{columns}
    \column{0.45\textwidth}
        \begin{itemize}
            \item<1-> \textbf{Dedutivo}
            \begin{itemize}
                \item Assumindo algumas premissas, a conclusão obrigatoriamente segue
            \end{itemize}
            \item<2-> \textbf{Não-dedutivo}
            \begin{itemize}
                \item<2-> \textbf{Abdução}
                \begin{itemize}
                    \item “Raciocínio do Detetive”
                \end{itemize}
                \item<3-> \textbf{Indução}
                \begin{itemize}
                    \item Dadas observações (limitadas em escopo e quantidade), podemos chegar a conclusões plausíveis
                \end{itemize}
            \end{itemize}
        \end{itemize}
        
    \column{0.55\columnwidth}
    \only<1>{
    \begin{block}{Exemplo de Dedução}
        \begin{itemize}
            \item Premissa 1: Todos os humanos são mortais
            \item Premissa 2: Sócrates é humano
            \item Conclusão: Sócrates é mortal
        \end{itemize}        
    \end{block}
    }
    \only<2>{
    \begin{block}{Exemplo de Abdução}
        \begin{itemize}
            \item Tem restos de frango assado pelo chão
            \item O cachorro está com restos de frango no pelo, especialmente ao redor da boca
            \item Logo, o cachorro comeu o frango assado
        \end{itemize}        
    \end{block}
    \begin{figure}
        \centering
        \includegraphics[width=0.4\linewidth]{slides/img/aula00/meme-hungrydog-abduction.jpg}
    \end{figure}
    }
    \only<3->{
    \begin{block}{Exemplo de Indução}
         Qual o próximo número na sequência abaixo?\\
         \alt<3>{$$1, 3, 5, 7, x$$}{$$1, 3, 5, 7, \textcolor{red}{217341}$$} \\    
        \only<4>{
            \scriptsize
            $$f(x) = \frac{18111}{2}x^4 - 90555x^3 + \frac{633885}{2}x^2 - 452773x + 217331$$
            }
    \end{block}
    \begin{figure}
        \centering
        \only<3>{
\includegraphics[width=0.35\linewidth]{slides/img/aula00/meme-happy-math.jpg}}
        \only<4>{\includegraphics[width=0.35\linewidth]{slides/img/aula00/meme-sad-math.jpg}}
    \end{figure}
    }
\end{columns}
\end{frame}

\section{Conceitos Básicos de \textit{Machine Learning}}

\begin{frame}{Desenvolvimento de uma aplicação com AM}
\only<1>{
    O projeto de um sistema que usa \textit{Machine Learning} pode ser descrito pela figura abaixo\footnote{\scriptsize Fluxos diferentes existem na literatura de \textit{Knowledge Discovery in Databases} e em Mineração de Dados.}
    }
\only<2>{
\textbf{Experimental Design}
    \begin{itemize}
        \item Modelar o problema como uma tarefa de \textit{Machine Learning}
        \item Como você definiria o problema de treinar um reconhecedor de rostos que é capaz de reconhecer 100 pessoas diferentes?
        \begin{itemize}
            \item Como um problema de \alert{classificação multiclasse} com 100 classes será muito difícil coletar uma base dados grande o bastante
            \item Como um problema de \alert{classificação binário} onde o sistema recebe um par de imagens é mais fácil coletar dados
        \end{itemize}
    \end{itemize}
}
\only<3>{
    \textbf{Data Collection}
    \begin{itemize}
        \item É possível coletar dados que sejam representativos do problema em questão?
        \item Existe algum dataset aberto que pode nos auxiliar?
        \item Quando a coleta de dados ou a anotação é cara, existem alternativas?
    \end{itemize}
}
\only<4>{
    \textbf{Dimensionality Reduction}
    \begin{itemize}
        \item Todos os atributos coletados são úteis para resolver a tarefa?
        \item Como selecionar apenas os atributos relacionados a variável alvo? 
        \item Existem formas de criar novos atributos que sejam melhores preditores para a variável alvo?
        \item Posso resumir meus atributos em um espaço menor?
    \end{itemize}
}
\only<5>{
    \textbf{Predictor Design}
    \begin{itemize}
        \item Existe algum modelo cujo viés indutivo seja aderente ao meu problema?
        \item Qual será a métrica utilizada para avaliação do desempenho do modelo?
        \item Qual meu \textit{budget} computacional?
        \item Qual a melhor hiperparametrização dos modelos selecionados?
    \end{itemize}
}
\only<6>{
    \textbf{Error Estimation}
    \begin{itemize}
        \item Meu modelo treinado generaliza?
        \item Qual é o erro esperado do meu modelo para novos dados?
        \item Coletar mais dados pode melhorar meus resultados?
    \end{itemize}

}

    \begin{figure}
        \centering
        \includegraphics[width=.5\textwidth]{slides/img/aula00/fluxo-ml.png}
        \caption{Representação simplificada do processo de \textit{Machine Learning}. Adaptado de \citeauthoryear{braga-neto_fundamentals_2020}}
    \end{figure}
    
\end{frame}


\begin{frame}{Tipos de Aprendizado}
    \begin{itemize}
        \item \textbf{Supervisionado}
        \begin{itemize}
            \item Dados de treino possuem rótulos
            \item \textit{Exemplo}: classificar fotos como gato ou cachorro a partir de 10 mil imagens rotuladas
        \end{itemize}
        \item \textbf{Não Supervisionado}
        \begin{itemize}
            \item Dados de treino não possuem rótulos
            \item \textit{Exemplo}: agrupar clientes em segmentos a partir do histórico de compras, sem rótulos
        \end{itemize}
        \item \textbf{Semi Supervisionado}
        \begin{itemize}
            \item Alguns dados de treino possuem rótulo outros não
            \item \textit{Exemplo}: detecção de fraude com poucas transações rotuladas e milhões sem rótulo
        \end{itemize}
        \item \textbf{Aprendizado por Reforço}
        \begin{itemize}
            \item Agente recebe “recompensa” devido a uma série de ações
            \item \textit{Exemplo}: AlphaGo aprende a jogar Go ganhando recompensa ao final da partida
        \end{itemize}
        \item \textbf{Auto Supervisionado}
        \begin{itemize}
            \item Parte das entradas servem como “supervisão”
            \item \textit{Exemplo}: pré-treino do BERT prevendo palavras mascaradas em uma frase
        \end{itemize}
    \end{itemize}
\end{frame}

\begin{frame}{Notação}
    Considerando um dataset formado por exemplos de uma relação $\big(\mlinstancei, \mltargeti\big )$
    \begin{itemize}
        \item $\mlinstance$ é nossa matriz de atributos $N\times M$ onde $N$ é o número de instâncias e $M$ é o número de atributos indexada por $\mlinstance^{(n)}_m$
        \item $\mltarget$ é nossa variável alvo $N\times 1$
        \item Objetivo é estimar $\mltargeti$ para uma instância desconhecida $\mlinstancei$
        \begin{itemize}
            \item Quando $\mltargeti$ é uma variável discreta: \textit{Problema de Classificação} 
            \item Quando $\mltargeti$ é uma variável contínua: \textit{Problema de Regressão}
        \end{itemize}
    \end{itemize}
\end{frame}


\begin{frame}{Viés Indutivo}
    \begin{columns}
        \begin{column}{0.25\textwidth}
            \begin{figure}
                \includegraphics[width=\textwidth]{slides/img/aula00/linear.png}
                \caption{Regressão Linear}
            \end{figure}
        \end{column}
        \begin{column}{0.25\textwidth}
            \begin{figure}
                \includegraphics[width=\textwidth]{slides/img/aula00/svm.png}
                \caption{SVM}
            \end{figure}
        \end{column}
        \begin{column}{0.25\textwidth}
            \begin{figure}
                \includegraphics[width=\textwidth]{slides/img/aula00/arvore.png}
                \caption{Árvore de Decisão}
            \end{figure}
        \end{column}
        \begin{column}{0.25\textwidth}
            \begin{figure}
                \includegraphics[width=\textwidth]{slides/img/aula00/randomforest.png}
                \caption{\textit{Random Forest}}
            \end{figure}
        \end{column}
    \end{columns}
\end{frame}


\begin{frame}{Efeito tesoura (\textit{Scissor Effect})}
    \begin{columns}
        \begin{column}{0.5\textwidth}
            \begin{itemize}
                \item Quando temos um conjunto de dados pequeno, regras de classificação mais simples tendem a obter erros de classificação menores
                \item Quando temos um conjunto de dados grande, regras de classificação mais complexas tendem a obter erros menores
            \end{itemize}
        \end{column}
        \begin{column}{0.5\textwidth}
            \begin{figure}
                \includegraphics[width=\textwidth]{slides/img/aula00/efeitotesoura.png}
                \caption{Adaptado de \citeauthoryear{braga-neto_fundamentals_2020}}
            \end{figure}
        \end{column}
    \end{columns}
\end{frame}

\begin{frame}{Efeito tesoura, exemplo numérico}
    \begin{columns}
        \begin{column}{0.55\textwidth}
            \centering
            \resizebox{\linewidth}{!}{%
            \begin{tikzpicture}[font=\footnotesize]
                \begin{axis}[
                    xlabel={$x$},
                    ylabel={$y$},
                    xmin=0, xmax=6.5,
                    ymin=-2, ymax=30,
                    xtick={1,2,3,4,5,6},
                    ymajorgrids=true, xmajorgrids=true,
                    grid style={dashed, gray!30},
                    legend pos=north west,
                    legend style={font=\scriptsize, fill=white, fill opacity=0.9, draw=gray!40},
                    legend cell align=left,
                    width=\linewidth, height=0.75\linewidth,
                    tick label style={font=\scriptsize},
                    label style={font=\small},
                ]
                    \addplot[color=red, line width=1pt, samples=120, domain=0.5:6.4, smooth]
                        {2.5 + (x-1) + (x-1)*(x-2) - (2/3)*(x-1)*(x-2)*(x-3) + (1/3)*(x-1)*(x-2)*(x-3)*(x-4)};
                    \addlegendentry{Polinômio grau 4}
                    \addplot[color=blue, line width=1pt, domain=0:6.4]
                        {2*x + 0.1};
                    \addlegendentry{Reta $\hat y = 2x + 0.1$}
                    \addplot[only marks, color=black, mark=*, mark size=2pt]
                        coordinates {(1,2.5)(2,3.5)(3,6.5)(4,7.5)(5,10.5)};
                    \addlegendentry{Treino}
                    \addplot[only marks, color=orange, mark=triangle*, mark size=4pt, mark options={draw=orange!70!black, line width=0.5pt}]
                        coordinates {(6,12)};
                    \addlegendentry{Teste $(6,\,12)$}
                \end{axis}
            \end{tikzpicture}}
            \\
            \scriptsize\textbf{Em $x=6$:} reta prediz $12.1$ (acerta), polinômio prediz $27.5$ (erra muito).
        \end{column}
        \begin{column}{0.45\textwidth}
            \begin{figure}
                \includegraphics[width=\textwidth]{slides/img/aula00/efeitotesoura.png}
                \caption{Adaptado de \citeauthoryear{braga-neto_fundamentals_2020}}
            \end{figure}
        \end{column}
    \end{columns}
\end{frame}

\begin{frame}{Trade-off Viés-Variância}
    \begin{columns}
        \begin{column}{0.55\textwidth}
            \begin{figure}
                \centering
                \includegraphics[width=\textwidth]{slides/img/aula00/tradeoffviesvar.png}
                \caption{Adaptado de \citeauthoryear{braga-neto_fundamentals_2020}.}
            \end{figure}
        \end{column}
        \begin{column}{0.45\textwidth}
            \begin{block}{Decomposição}
                \small
                \[
                    \text{Erro}_{\text{teste}}(x) = \text{Viés}^{2}(x) + \text{Variância}(x) + \sigma^{2}
                \]
            \end{block}
            \begin{itemize}
                \item \textbf{Viés alto}: modelo simples demais, não captura o padrão.
                \item \textbf{Variância alta}: modelo sensível à amostra de treino.
                \item \textbf{Ruído} ($\sigma^{2}$): limite irredutível.
            \end{itemize}
            \footnotesize Derivação completa em \textit{Regularização} (Aula 08).
        \end{column}
    \end{columns}
\end{frame}


\section{Modelos Lineares}

\begin{frame}{Modelos Lineares}
\begin{itemize}
    \item Modelos lineares baseados em variações dos Método dos Mínimos Quadrados talvez sejam o tipo de modelo mais estudado por todas as áreas do conhecimento.
    \item Faremos uma breve revisão com viés de \textit{Machine Learning}
\end{itemize}

\end{frame}
\subsection{Regressão Linear}

\begin{frame}{Intuição}
\only<1-2>{
Queremos ajustar um \textit{hiperplano}\footnote{No $\mathbb{R}^2$ um hiperplano é equivalente a uma reta} aos nossos dados

$$\hat{\mltarget}=\mlinstance_1\cdot\theta_1+\theta_0$$
}
\only<3-4>{
Qual é a melhor \alert{hipótese}?
\[J_{\text{SSR}}=\sum_{i=1}^{N}\left(\mltargethati - \mltargeti\right)^2\]
}
\only<2->{
\begin{columns}[onlytextwidth]
    \begin{column}{0.1\textwidth}
        \scriptsize
        \begin{table}[]
            \centering
            \begin{tabular}{|c|c|}
            \hline
                $\mlinstance$ & $\mltarget$  \\
                \hline
                 1 & 20  \\
                 \hline
                 2 & 10 \\
                 \hline
                 3 & 30 \\
                 \hline
            \end{tabular}
            \caption{\scriptsize Conjunto de Dados}
        \end{table}
    \end{column}
    \begin{column}{0.4\textwidth}
        \begin{tikzpicture}[font=\tiny]
        \begin{axis}[
            %scale only axis,
            width=\textwidth,
            legend pos=south east,% <-- this line controls the width
            xlabel={$\mlinstance_1$},
            ylabel={$\mltarget$},
            xmin=0, xmax=5,
            ymin=0, ymax=50,
            xtick={0,1,2,3,4,5},
            ytick={0,10,20,30,40,50},
            ymajorgrids=true,
            xmajorgrids=true,
            grid style=dashed,
        ]
        
        \addplot[
            only marks,
            color=red,
            mark=*,
            ]
            coordinates {
            (1,20)(2,10)(3,30)
            };
        \addplot[
            color=blue,
            line width=1pt,
            ]
            coordinates {
            (0,15)(5,40)
            };
        \addplot[
            color=green,
            line width=1pt,
            ]
            coordinates {
            (0,20)(5,20)
            };
        \node at (4,40) {\tiny Hipótese 2};
        \node at (4,22) {\tiny Hipótese 1};
        \end{axis}
        \end{tikzpicture}
        
        
    \end{column}
    \begin{column}{0.5\textwidth}
        \begin{block}{\only<4>{\emoji{check-mark-button}} Hipótese 1}
        \scriptsize
            \[ \theta_0=20, \theta_1=0\]
            \[\hat{\mltarget}=0\mlinstance_1+20\]
            \only<4>{$J_{\text{SSR}}=(20-20)^2+(10-20)^2+(30-20)^2=200$}
        \end{block}
        \begin{block}{\only<4>{\emoji{cross-mark}} Hipótese 2}
        \scriptsize
            \[ \theta_0=15, \theta_1=5\]
            \[\hat{\mltarget}=5\mlinstance_1+15\]
            \only<4>{$J_{\text{SSR}}=(20-20)^2+(10-25)^2+(30-30)^2=225
$}
        \end{block}
    \end{column}
\end{columns}
}   
\end{frame}
\normalsize
\begin{frame}{Treinamento}
Como encontrar a melhor hipótese?
\begin{itemize}
    \item\(\argmin_{\theta_0,\theta_1}\left[J_{\text{SSR}}\right]\)
    \item \(J_{\text{SSR}}=\sum_{i=1}^{N}\left(\mltargethati - \mltargeti\right)^2=(\mltargethat - \mltarget)^\top(\mltargethat - \mltarget)\)
    \item \(\hat{\mltarget}=\mlinstance_1\cdot\theta_1+\theta_0\)
\end{itemize}

    \[
    \begin{array}{|c|}
      \hline
        \mlinstance_1  \\
        \hline
        1   \\
        \hline
        2   \\
        \hline
        3   \\
        \hline
      \end{array}
      \quad
      \begin{array}{|c|}
      \hline
        \mltarget  \\
        \hline
        20   \\
        \hline
        10   \\
        \hline
        30   \\
        \hline
      \end{array}
      \quad
      \begin{array}{|c|}
      \hline
        \theta_0  \\
        \hline
      \end{array}
      \quad
      \begin{array}{|c|}
      \hline
        \theta_1  \\
        \hline
      \end{array}
    \]
\end{frame}

\begin{frame}{\textit{Notation Trick}}

\begin{itemize}
    \item Colocamos o parâmetro livre junto do vetor $\theta$ fazendo um \textit{truque de notação}
    \item Imagine que existe um novo atributo cujo valor é 1 para todas as instâncias
    \item Assim todos os parâmetros treináveis podem ser tratados da mesma forma.
    \item $\mllinearregression$
\end{itemize}


    \[
    \begin{array}{|>{\columncolor{PastelPink}}c|>{\columncolor{PastelGreen}}c|}
      \hline
        \mlinstance_0 &  \mlinstance_1  \\
        \hline
        1 & 1   \\
        \hline
        1 & 2   \\
        \hline
        1 & 3   \\
        \hline
      \end{array}_{N\times M}
      \quad
      \begin{array}{|c|}
      \hline
        \mltarget  \\
        \hline
        20   \\
        \hline
        10   \\
        \hline
        30   \\
        \hline
      \end{array}_{N\times 1}
      \quad
      \begin{array}{|c|}
      \hline
      \bm{{\theta}} \\
        \hline
        \cellcolor{PastelPink}\theta_0  \\
        \hline
        \cellcolor{PastelGreen}\theta_1  \\
        \hline
      \end{array}_{M\times 1}
    \]
\end{frame}

\begin{frame}{Treinamento}
Como encontrar a melhor hipótese?
\begin{itemize}
    \item\(\argmin_{\theta_0,\theta_1}\left[J_{\text{SSR}}\right]\)
    \item \(J_{\text{SSR}}=\sum_{i=1}^{N}\left(\mltargethati - \mltargeti\right)^2=(\mltargethat - \mltarget)^\top(\mltargethat - \mltarget)\)
    \item \(\mllinearregression\)
\end{itemize}
\begin{columns}
    \begin{column}{.5\textwidth}
    \[
    \underset{N\times M}{
    \begin{array}{|c|c|}
      \hline
        \mlinstance_0 &  \mlinstance_1  \\
        \hline
        1 & 1   \\
        \hline
        1 & 2   \\
        \hline
        1 & 3   \\
        \hline
      \end{array}
      }
      \quad
      \underset{N\times 1}{
      \begin{array}{|c|}
      \hline
        \mltarget  \\
        \hline
        20   \\
        \hline
        10   \\
        \hline
        30   \\
        \hline
      \end{array}}
      \quad
      \underset{M\times 1}{
      \begin{array}{|c|}
      \hline
      \bm{{\theta}} \\
        \hline
        \theta_0  \\
        \hline
        \theta_1  \\
        \hline
      \end{array}}
    \]    
    \end{column}
    \begin{column}{.5\textwidth} 
        \begin{block}{Sabendo que..}
            \begin{itemize}
                \item A soma dos quadrados dos resíduos é uma \textit{função de custo convexa} para o problema da regressão linear
                \item Logo, \textit{possui apenas um mínimo}
                \item Podemos utilizar o gradiente da função de custo em um processo iterativo de otimização
            \end{itemize}
        \end{block}

    \end{column}
\end{columns}
\end{frame}

\begin{frame}{Treinamento}
    
        Vamos derivar $J_{\text{SSR}}$ em relação a $\theta$ e igualar a 0
        \begin{align*}
            \frac{\partial J_{\text{SSR}}}{\partial\theta}&=\frac{\partial ((\mltargethat - \mltarget)^\top(\mltargethat - \mltarget))}{\partial\theta}\\
            &=\frac{\partial (\mltargethat^\top\mltargethat - \mltargethat^\top\mltarget-\mltarget^\top\mltargethat + \mltarget^\top\mltarget)}{\partial\theta}\\
            &=\frac{\partial (\mltargethat^\top\mltargethat - 2\mltargethat^\top\mltarget + \mltarget^\top\mltarget)}{\partial\theta}\\
            &=\frac{\partial (\bm{\theta}^\top\mlinstance^\top\mlinstance\bm{\theta} - 2\bm{\theta}^\top\mlinstance^\top\mltarget + \mltarget^\top\mltarget)}{\partial\theta}\\
            &=\frac{\partial (\bm{\theta}^\top\mlinstance^\top\mlinstance\bm{\theta})}{\partial\theta} - \frac{\partial(2\bm{\theta}^\top\mlinstance^\top\mltarget)}{\partial\theta}\\
            &=2\mlinstance^\top\mlinstance\bm{\theta} - 2\mlinstance^\top\mltarget\\
            &=2\mlinstance^\top(\mlinstance\bm{\theta} - \mltarget)
        \end{align*}
\end{frame}

\begin{frame}{Treinamento}
    
        Vamos derivar $J_{\text{SSR}}$ em relação a $\theta$ e igualar a 0
        \begin{align*}
            2\mlinstance^\top(\mlinstance\bm{\theta} - \mltarget) &= 0\\
            \mlinstance^\top(\mlinstance\bm{\theta} - \mltarget) &= 0\\
            \mlinstance^\top\mlinstance\bm{\theta} - \mlinstance^\top\mltarget &= 0\\
            \mlinstance^\top\mlinstance\bm{\theta} &= \mlinstance^\top\mltarget\\
            \bm{\theta} &= (\mlinstance^\top\mlinstance)^{-1}\mlinstance^\top\mltarget\\
        \end{align*}
\end{frame}

\begin{frame}{Solução analítica para $\bm{\theta}$}

        $$\bm{\theta} = (\mlinstance^\top\mlinstance)^{-1}\mlinstance^\top\mltarget$$
\only<1>{
     \[
     \left (\mlexampleinstancetranspose
    \times
        \mlexampleinstance
      \right)^{-1}
      \times
      \mlexampleinstancetranspose
      \times
      \mlexampletarget
      =
      \underset{M\times 1}{
      \bm{{\theta}}}
    \] } 

    \only<2>{
     \[
     \left (
     \underset{M\times M}{
     \begin{array}{|c|c|}
        \hline
        3 & 6  \\
        \hline
        6 & 14  \\
        \hline
      \end{array}}
      \right)^{-1}
      \times
      \mlexampleinstancetranspose
      \times
      \mlexampletarget
      =
      \underset{M\times 1}{
      \bm{{\theta}}}
    \] }

    \only<3>{
     \[
     \underset{M\times M}{
     \begin{array}{|c|c|}
        \hline
        2.33 & -1  \\
        \hline
        -1 & 0.5  \\
        \hline
      \end{array}}
      \times
      \mlexampleinstancetranspose
      \times
      \mlexampletarget
      =
      \underset{M\times 1}{
      \bm{{\theta}}}
    \] }

    \only<4>{
     \[
     \underset{M\times N}{
     \begin{array}{|c|c|c|}
        \hline
        1.33 & 0.333  & -0.666\\
        \hline
        -0.5 & 0 & 0.5\\
        \hline
      \end{array}}
      \times
      \mlexampletarget
      =
    \underset{M\times 1}{
      \bm{{\theta}}}
    \] }

    \only<5>{
     \[
     \underset{M\times 1}{
     \begin{array}{|c|}
        \hline
        10\\
        \hline
        5\\
        \hline
      \end{array}}
      =
      \underset{M\times 1}{
      \bm{{\theta}}}
    \] 
    \begin{alertblock}{Sobre a solução analítica para $\bm{\theta}$}
    \begin{itemize}
        \item Esse é o vetor de pesos que minimiza a soma dos quadrados dos resíduos ($J_{\text{SSR}}$).
        \item Repare que existe uma operação de inversão da matriz $\mlinstance^\top\mlinstance$. Essa operação gera um custo computacional $\mathcal{O}(M^3)$ em relação à quantidade de atributos do problema
        \begin{itemize}
            \item Esse custo pode inviabilizar a solução analítica
        \end{itemize}
    \end{itemize}
    
    \end{alertblock}
    } 
\end{frame}

\begin{frame}{Qual é a melhor \alert{hipótese}?}
    \begin{columns}[onlytextwidth]
    \begin{column}{0.1\textwidth}
        \scriptsize
        \begin{table}[]
            \centering
            \begin{tabular}{|c|c|}
            \hline
                $\mlinstance$ & $\mltarget$  \\
                \hline
                 1 & 20  \\
                 \hline
                 2 & 10 \\
                 \hline
                 3 & 30 \\
                 \hline
            \end{tabular}
            \caption{\scriptsize Conjunto de Dados}
        \end{table}
    \end{column}
    \begin{column}{0.4\textwidth}
        \begin{tikzpicture}[font=\tiny]
        \begin{axis}[
            %scale only axis,
            width=\textwidth,
            legend pos=south east,% <-- this line controls the width
            xlabel={$\mlinstance_1$},
            ylabel={$\mltarget$},
            xmin=0, xmax=5,
            ymin=0, ymax=50,
            xtick={0,1,2,3,4,5},
            ytick={0,10,20,30,40,50},
            ymajorgrids=true,
            xmajorgrids=true,
            grid style=dashed,
        ]
        
        \addplot[
            only marks,
            color=red,
            mark=*,
            ]
            coordinates {
            (1,20)(2,10)(3,30)
            };
        \addplot[
            color=blue,
            line width=1pt,
            ]
            coordinates {
            (0,15)(5,40)
            };
        \addplot[
            color=green,
            line width=1pt,
            ]
            coordinates {
            (0,20)(5,20)
            };
        \addplot[
            color=magenta,
            line width=1pt,
            ]
            coordinates {
            (0,10)(5,35)
            };
        \node at (4,40) {\tiny Hipótese 2};
        \node at (4,25) {\tiny Hipótese 3};
        \node at (4,17) {\tiny Hipótese 1};
        \end{axis}
        \end{tikzpicture}
        
    \end{column}
    \begin{column}{0.5\textwidth}
        \begin{block}{\emoji{cross-mark} Hipótese 1}
        \scriptsize
            \[ \theta_0=20, \theta_1=0, \hat{\mltarget}=0\mlinstance_1+20\]
            $J_{\text{SSR}}=(20-20)^2+(10-20)^2+(30-20)^2=200$
        \end{block}
        \begin{block}{\emoji{cross-mark} Hipótese 2}
        \scriptsize
            \[ \theta_0=15, \theta_1=5, \hat{\mltarget}=5\mlinstance_1+15\]
            $J_{\text{SSR}}=(20-20)^2+(10-25)^2+(30-30)^2=225$
        \end{block}
        \begin{block}{\emoji{check-mark-button} Hipótese 3}
        \scriptsize
            \[ \theta_0=10, \theta_1=5,\hat{\mltarget}=5\mlinstance_1+10\]
            $J_{\text{SSR}}=(20-15)^2+(10-20)^2+(30-25)^2=150$
        \end{block}
    \end{column}
\end{columns}
\end{frame}

\begin{frame}{Descida de Gradiente}
\begin{itemize}
    \item Algoritmo Iterativo
    \item $\bm{\theta}_{t+1}=\bm{\theta}_{t}-\eta\nabla_ J$
    \item $\bm{\theta}_{0}$ inicializado com valores aleatórios
    \item $\eta\in(0,1]$ taxa de aprendizado
    \item Valores típicos para $\eta$ são $10^{-3}, 10^{-4}$ ou $10^{-5}$
\end{itemize}    
\end{frame}

\begin{frame}{Descida de Gradiente}
\begin{columns}
    \begin{column}{.45\linewidth}
    \begin{itemize}
        \item Ao lado temos um algoritmo para \textit{Batch Gradient Descent} 
        \item Se selecionarmos apenas uma instância $\mlinstancei$ para estimar o valor do gradiente teremos o \textit{Stochastic Gradient Descent}
        \item A diferença entre $J_{\text{SSR}}$ de diferentes iterações pode ser usado como condição de convergência
    \end{itemize}

    \end{column}
    \begin{column}{.55\linewidth}
        \begin{algorithmic}[1]
        \Require $\mlinstance_{n\times m}$ matriz de instâncias
        \Require $\mltarget_{n\times 1}$ variável alvo
        \Require $\eta\in(0,1]$ taxa de aprendizado
        \Procedure{Gradient Descent}{}
        \State $\bm{\theta} \sim \mathcal{N}(0,1)$ 
        \Repeat 
            \State $\bm{\theta} \gets \bm{\theta}-\eta\mlinstance^\top(\mlinstance\bm{\theta}-\mltarget)$
            \State $\mltargethat\gets\mlinstance\bm{\theta}$
            \State $J_{\text{SSR}} \gets (\mltarget-\mltargethat)^\top(\mltarget-\mltargethat)$ 
        \Until{Condição de convergência}
        \State \Return $\bm{\theta}$
        \EndProcedure
        \end{algorithmic}        
    \end{column}
\end{columns}
    
\end{frame}


\begin{frame}{Descida de Gradiente}
\begin{table}[]
    \centering
    \begin{tabular}{cc}
        \toprule
       Descida de Gradiente  & Equação Normal \\
       \hline
       $\bm{\theta}_{t+1}=\bm{\theta}_{t}-\eta\mlinstance^\top(\mlinstance\bm{\theta}_{t}-\mltarget)$&$\bm{\theta} = (\mlinstance^\top\mlinstance)^{-1}\mlinstance^\top\mltarget$\\
        solução iterativa & solução analítica \\
          Possui o hiperparâmetro $\eta$ & Sem hiperparâmetros \\
     arbitrariamente próxima ao melhor $\bm{\theta}$  & encontra exatamente o melhor $\bm{\theta}$ \\
          $\mathcal{O}(i\cdot N\cdot M)$ em implementações vetoriais& $\mathcal{O}(M^3)$ \\
        \bottomrule
    \end{tabular}
\end{table}
\end{frame}


\subsection{Regressão Logística}
\begin{frame}{Como aplicar Regressão Linear para Classificação?}
A \textit{imagem} da Regressão Linear é incompatível com uma classificação
\begin{itemize}
    \item $p\left(\mltargethati=c|\mlinstancei\right)\in[0,1]$
    \item $\mlinstancei\bm{\theta}\in(-\infty,+\infty)$
\end{itemize}

Precisamos encontrar algum valor relacionado a probabilidade de pertencimento a uma classe $p\left(\mltargethati=c|\mlinstancei\right)$ que tenha um contradomínio similar a Regressão Linear $\mlinstancei\bm{\theta}$
    
\end{frame}

\begin{frame}{Como aplicar Regressão Linear para Classificação?}

\begin{itemize}
    \item<1-> $p\left(\mltargethati=c|\mlinstancei\right)\in[0,1]$: probabilidade de pertencimento a classe $c$ 
    \item<2-> $\frac{p\left(\mltargethati=c|\mlinstancei\right)}{1-p\left(\mltargethati=c|\mlinstancei\right)}\in[0,+\infty)$: chance de pertencimento a classe $c$
    \item<3-> $\ln{\left(\frac{p\left(\mltargethati=c|\mlinstancei\right)}{1-p\left(\mltargethati=c|\mlinstancei\right)}\right)}\in(-\infty,+\infty)$: log da chance de pertencimento a classe $c$: 
\end{itemize}

\end{frame}

\begin{frame}{Aplicando Regressão Linear para estimar o log das chances}
\small
\begin{align*}
    \ln{\left(\frac{p\left(\mltargethati=c|\mlinstancei\right)}{1-p\left(\mltargethati=c|\mlinstancei\right)}\right)} &= \mlinstancei\bm{\theta}\\
    \frac{p\left(\mltargethati=c|\mlinstancei\right)}{1-p\left(\mltargethati=c|\mlinstancei\right)} &= e^{\mlinstancei\bm{\theta}}\\
    p\left(\mltargethati=c|\mlinstancei\right) &= e^{\mlinstancei\bm{\theta}}\left(1-p\left(\mltargethati=c|\mlinstancei\right)\right)\\
    p\left(\mltargethati=c|\mlinstancei\right) +p\left(\mltargethati=c|\mlinstancei\right)e^{\mlinstancei\bm{\theta}} &= e^{\mlinstancei\bm{\theta}}\\
    p\left(\mltargethati=c|\mlinstancei\right)(1+e^{\mlinstancei\bm{\theta}}) &= e^{\mlinstancei\bm{\theta}}\\
    p\left(\mltargethati=c|\mlinstancei\right) &= \frac{e^{\mlinstancei\bm{\theta}}}{(1+e^{\mlinstancei\bm{\theta}})}= \frac{1}{e^{-\mlinstancei\bm{\theta}}(1+e^{\mlinstancei\bm{\theta}})}\\
    p\left(\mltargethati=c|\mlinstancei\right) &= \frac{1}{\left(1+e^{-\mlinstancei\bm{\theta}}\right)}\\
\end{align*}
\end{frame}

\begin{frame}{Regressão Logística}
\begin{itemize}
    \item Ao adaptar a Regressão Linear para a tarefa de classificação binária obtemos a Regressão Logística
    \item $\mltargethat = \frac{1}{\left(1+e^{-\mlinstancei\bm{\theta}}\right)}$
    \item Precisamos de uma nova função de custo
    \begin{itemize}
        \item $J_{\text{SSR}}$ não é convexo para a Regressão Logística
        \item $J_{\text{BCE}}=-\frac{1}{N}\sum_{i=1}^{N}\Big[\mltargeti\ln\mltargethati +(1-\mltargeti)\ln{(1-\mltargethati)}\Big]$
    \end{itemize}
\end{itemize}
    
\end{frame}

\begin{frame}{Treinamento}
Vamos derivar $J_{\text{BCE}}$ em relação aos $\bm{\theta}$
\begin{align*}
    \frac{\partial J_{\text{BCE}}}{\partial\bm{\theta}} &=     \frac{\partial (-\mltarget\ln\mltargethat -(1-\mltarget)\ln{(1-\mltargethat)})}{\partial\bm{\theta}}\\
    &= -\frac{\partial (\mltarget\ln\mltargethat)}{\partial\bm{\theta}} -\frac{\partial((1-\mltarget)\ln{(1-\mltargethat)})}{\partial\bm{\theta}}\\
    &= -\mltarget\frac{\partial (\ln\mltargethat)}{\partial\bm{\theta}} -(1-\mltarget)\frac{\partial\ln{(1-\mltargethat)}}{\partial\bm{\theta}}\\
    &= -\mltarget\frac{\partial (\ln\mltargethat)}{\partial\mltargethat}\frac{\partial\mltargethat}{\partial\bm{\theta}} -(1-\mltarget)\frac{\partial\ln{(1-\mltargethat)}}{\partial\mltargethat}\frac{\partial\mltargethat}{\partial\bm{\theta}}\\
    &= -\mlinstance^\top\mltarget\frac{1}{\mltargethat}\mltargethat(1-\mltargethat) -\mlinstance^\top(1-\mltarget)\left(-\frac{1}{1-\mltargethat}\right)\mltargethat(1-\mltargethat)\\
    &= \mlinstance^\top\left(-\mltarget+\mltarget\mltargethat +\mltargethat-\mltargethat\mltarget\right)\\
    &= \mlinstance^\top\left(\mltargethat-\mltarget\right)\\
\end{align*}
    
\end{frame}

\begin{frame}{Regressão Logística}
\begin{itemize}
    \item Função de Custo $J_{\text{BCE}}$ é convexa, mas não admite solução analítica
    \item O gradiente tem exatamente a mesma forma da Regressão Linear
    \item Gera uma fronteira de decisão linear (hiperplano) no espaço de atributos
    \end{itemize}
    
\end{frame}

\section*{Referências e Leituras}
%------------------------------------------------
\begin{frame}{Leituras Recomendadas}
    \leiturasize
    \begin{itemize}
        \item Seções 3.1 a 3.3 de: \fullcite{abu-mostafa_learning_2012}
    \end{itemize}

\end{frame}


%------------------------------------------------

\begin{frame}{Referências}
    \printbibliography

\end{frame}

%------------------------------------------------

% \begin{frame}
%     \Huge{\centerline{\textbf{The End}}}
% \end{frame}

%----------------------------------------------------------------------------------------

\end{document}